% Fichier complété par tools/complete_missing_corrections.py.
% Source : TEC/TEC-04-Meq-Jeq/34_ControlX/34_ControlX.tex
% Statut : rédigé (3 question(s)).
\begin{corrigebox}[Corrigé rédigé]
\CorrigeQuestion{1}
La couronne est fixe. Pour le train épicycloïdal simple,
            la relation de Willis donne
            $(\omega_1-\omega_3)/(\omega_0-\omega_3)=-R_b/R_m$ avec $\omega_0=0$.
            D'où $\omega_3=\omega_1R_m/(R_m+R_b)$. La courroie ne glissant pas sur la
            poulie de rayon $R_p$, $v=R_p\omega_3$, soit
            \[\boxed{v=R_p\frac{R_m}{R_m+R_b}\,\omega_{1/0}}.\]
\CorrigeQuestion{2}
$J_{eq,m}=J_m+\sum J_i(\omega_i/\omega_m)^2+M(v/\omega_m)^2$.
                Avec $v/\omega_m=R_pR_m/(R_m+R_b)$, la contribution du chariot est
                $M[R_pR_m/(R_m+R_b)]^2$.
\CorrigeQuestion{3}
Ramenée à la translation, la masse équivalente vaut
                $M_{eq}=M+\sum J_i(\omega_i/v)^2$. En particulier l'inertie moteur ajoute
                $J_m[(R_m+R_b)/(R_pR_m)]^2$.
\end{corrigebox}
